% ── Preprint template — swap style by changing the \usepackage line ───────────
% Available styles: styles/seed
\documentclass[11pt]{article}
\usepackage{styles/seed}

% ── Paper metadata ────────────────────────────────────────────────────────────
\title{How we found AGI by using one A100 node}

\author{
  Patrick Haller\textsuperscript{*} Filipe Laitenberger\textsuperscript{*}
}

\affiliations{Humboldt-Universität zu Berlin}

\authornotes{%
  \textsuperscript{*}Corresponding author%
}

% Optional: arXiv watermark on left margin
% \arxivwatermark{2506.XXXXX}{cs.LG}

% Logo shown in bottom-right of abstract box (only when fill is off)
\preprintlogo{assets/hu-logo.png}

% Abstract style toggle:
%   \abstractfilledtrue   — pale blue fill, no logo (default)
%   \abstractfilledfalse  — white fill + logo in bottom-right corner
\abstractfilledfalse

% Abstract metadata
\preprintdate{June 3, 2026}
\correspondence{patrick.haller@hu-berlin.de}

% ─────────────────────────────────────────────────────────────────────────────
\begin{document}

\maketitle

\begin{abstract}
Linear recurrent models have recently emerged as competitive alternatives to
Transformers, offering sub-quadratic training complexity and constant-memory
inference. Among these, \textbf{DeltaProduct}~\cite{vaswani2017attention}
demonstrated that replacing the scalar forgetting factor of linear attention
with a product of rank-1 updates significantly improves state-tracking and
associative recall. However, the original formulation applies a single global
gate per time step, leaving fine-grained, per-channel dynamics unexploited.

We introduce \textbf{DeltaProduct2}, which augments DeltaProduct with
\emph{fine-grained input, forget, and output gates} that operate independently
on each hidden dimension. This allows the model to selectively retain or
discard information at the feature level, mirroring the expressivity gains seen
in gated RNNs without sacrificing the hardware-efficient delta-rule update.
On synthetic state-tracking benchmarks (Flip-Flop, Parity, Dyck languages)
DeltaProduct2 matches or outperforms prior linear recurrences while training
$1.4\times$ faster than comparable chunk-parallel baselines. Language-modelling
experiments on the Pile at the 360M and 1.3B parameter scales show consistent
perplexity improvements over both DeltaNet and vanilla DeltaProduct,
approaching Transformer quality with a fraction of the inference cost.
\end{abstract}

% ─────────────────────────────────────────────────────────────────────────────
\section{Introduction}

The dominant sequence-modelling paradigm remains the Transformer~\cite{vaswani2017attention},
whose self-attention mechanism achieves strong empirical results but scales
quadratically in both compute and memory with sequence length.
This has spurred a renewed interest in linear recurrent models that retain
the expressive power of attention while offering $\mathcal{O}(L)$ inference.

\textbf{DeltaNet}~\cite{vaswani2017attention} and its successor
\textbf{DeltaProduct} extend linear attention with a delta-rule memory update,
achieving markedly better state-tracking---the ability to maintain and
manipulate discrete state across long contexts.
Despite these advances, a key limitation remains: the forgetting mechanism
applies a single scalar (or low-rank) gate shared across all hidden dimensions,
which constrains the model's capacity to independently track multiple
concurrent states.

We address this with \textbf{DeltaProduct2}.
Our contributions are:
\begin{itemize}
  \item A fine-grained gating scheme for the DeltaProduct recurrence that
        decouples forget and input signals per hidden dimension (\S\ref{sec:method}).
  \item A hardware-efficient chunk-parallel implementation compatible with
        modern accelerators (\S\ref{sec:method}).
  \item Systematic evaluation on state-tracking benchmarks and
        language-modelling at scale (\S\ref{sec:experiments}).
\end{itemize}

\subsection{Related Work}

\paragraph{Linear recurrent models.}
Linear attention~\cite{vaswani2017attention}, S4, Mamba, and RWKV all forgo
full softmax attention in favour of recurrent state updates.
We situate DeltaProduct2 within this landscape and contrast it with the
closest competitors.

\paragraph{Gating in recurrent networks.}
LSTMs and GRUs introduced per-dimension gates to control information flow.
Recent work has revisited this idea in the linear-attention setting (e.g.,
GLA, RetNet with gating variants).
DeltaProduct2 brings fine-grained gating specifically to the delta-rule
update, a combination not explored in prior work.

\paragraph{State tracking.}
\ldots [placeholder: discuss Flip-Flop, Parity, formal languages, prior
linear-attention results]

% ─────────────────────────────────────────────────────────────────────────────
\section{Background}
\label{sec:background}

\subsection{Linear Attention as Associative Memory}

Linear attention can be viewed as a key-value associative memory
$\mathbf{S}_t \in \mathbb{R}^{d_v \times d_k}$ updated at each step:
\begin{equation}
  \mathbf{S}_t = \mathbf{S}_{t-1} + \mathbf{v}_t \mathbf{k}_t^\top, \qquad
  \mathbf{o}_t = \mathbf{S}_t \mathbf{q}_t.
\end{equation}
The absence of a forgetting term causes unbounded memory growth and
poor state-tracking on tasks requiring selective overwriting.

\subsection{DeltaNet and DeltaProduct}

DeltaNet~\cite{vaswani2017attention} introduces a delta-rule correction:
\begin{equation}
  \mathbf{S}_t = \mathbf{S}_{t-1}
    + \beta_t \bigl(\mathbf{v}_t - \mathbf{S}_{t-1}\mathbf{k}_t\bigr)\mathbf{k}_t^\top,
\end{equation}
where $\beta_t \in (0,1)$ is a scalar gate.
DeltaProduct generalises this by replacing the single update with a
\emph{product} of rank-1 corrections, improving memory capacity and
state-tracking at the cost of a more complex recurrence.

% ─────────────────────────────────────────────────────────────────────────────
\section{DeltaProduct2}
\label{sec:method}

\subsection{Fine-Grained Gating}

We parameterise separate input gate $\mathbf{i}_t \in \mathbb{R}^{d_v}$,
forget gate $\mathbf{f}_t \in \mathbb{R}^{d_v}$, and output gate
$\mathbf{o}_t \in \mathbb{R}^{d_v}$, each produced by a small linear
projection of the input $\mathbf{x}_t$:
\begin{align}
  \mathbf{i}_t &= \sigma(W_i \mathbf{x}_t + b_i), \\
  \mathbf{f}_t &= \sigma(W_f \mathbf{x}_t + b_f), \\
  \mathbf{g}_t &= \sigma(W_g \mathbf{x}_t + b_g).
\end{align}
The recurrence becomes:
\begin{equation}
  \mathbf{S}_t = \mathrm{diag}(\mathbf{f}_t)\,\mathbf{S}_{t-1}
    + \mathrm{diag}(\mathbf{i}_t)\,
      \bigl(\mathbf{v}_t - \mathbf{S}_{t-1}\mathbf{k}_t\bigr)\mathbf{k}_t^\top.
\end{equation}

\subsection{Hardware-Efficient Implementation}

\ldots [placeholder: describe chunk-parallel scan, tiling strategy,
memory layout, FLOP count comparison]

\subsection{Connection to Prior Models}

Setting $\mathbf{i}_t = \mathbf{f}_t = \beta_t \mathbf{1}$ recovers
vanilla DeltaNet, and setting $\mathbf{f}_t = \mathbf{1}$ recovers
standard linear attention---showing DeltaProduct2 strictly generalises
both.

% ─────────────────────────────────────────────────────────────────────────────
\section{Experiments}
\label{sec:experiments}

\subsection{State-Tracking Benchmarks}

We evaluate on three synthetic tasks designed to probe discrete state
manipulation: \textbf{Flip-Flop}, \textbf{Parity}, and
\textbf{Dyck-$k$} bracket matching.
\ldots [placeholder: table of results]

\subsection{Language Modelling}

We train models at 360M and 1.3B parameters on the Pile corpus and
report perplexity on the held-out validation set.
\ldots [placeholder: perplexity table, training curve figure]

\subsection{Ablations}

\paragraph{Gate granularity.}
We compare scalar, per-head, and per-dimension gates to isolate the
contribution of fine-grained gating.

\paragraph{Number of delta steps.}
Following DeltaProduct, we sweep the number of rank-1 update steps
$K \in \{1, 2, 4, 8\}$ with and without fine-grained gates.

% ─────────────────────────────────────────────────────────────────────────────
\section{Conclusion}

DeltaProduct2 extends the DeltaProduct recurrence with fine-grained,
per-dimension gates that independently control information flow into and
out of the associative memory.
The result is a linear recurrent model that achieves stronger
state-tracking than its predecessors while remaining computationally
efficient.
Future work includes scaling beyond 1.3B parameters, exploring
multi-head variants with shared gate banks, and applying DeltaProduct2
to long-document retrieval tasks where fine-grained forgetting is
especially beneficial.

% ── Bibliography ──────────────────────────────────────────────────────────────
\bibliographystyle{plainnat}
\bibliography{references}

\end{document}
